% Figure 4. A cut annulus and the overlapping Hartogs set P.
\begin{tikzpicture}[field/diagram, scale=1.03]
    \def\CutAnnulus{
        ({2.5*cos(-78)},{1.8+1.6*sin(-78)})
        arc[start angle=-78, end angle=258, x radius=2.5, y radius=1.6]
        -- ({2.05*cos(258)},{1.8+1.22*sin(258)})
        arc[start angle=258, end angle=-78, x radius=2.05, y radius=1.22]
        -- cycle
    }
    \def\HartogsBoundary{
        (-1.03,1.3) -- (1.03,1.3) -- (1.03,0.66)
        -- (0.37,0.66) -- (0.37,-1.0) -- (1.03,-1.0)
        -- (1.03,-1.65) -- (-1.03,-1.65) -- (-1.03,-1.0)
        -- (-0.37,-1.0) -- (-0.37,0.66) -- (-1.03,0.66) -- cycle
    }
    \fill[FieldRegion] \CutAnnulus;
    \fill[FieldRegion] \HartogsBoundary;
    % Darker gray records the actual overlap, not an additional set.
    \begin{scope}
        \clip \CutAnnulus;
        \fill[FieldEmphasis] \HartogsBoundary;
    \end{scope}
    \draw[field/boundary] \CutAnnulus;
    \draw[field/boundary] \HartogsBoundary;
    \draw[field/hidden] (-1.03,0.66) -- (-1.03,-1.0)
                        (1.03,0.66) -- (1.03,-1.0);
    \node[inner sep=0.5pt] at (0,3.20) {$A$};
    \node at (0,0.25) {$P$};
\end{tikzpicture}
